\section{Open Questions}

The following five questions remain genuinely open after this replication.
Each is unresolved by both the paper and this reimplementation, and each
admits a concrete follow-up study using free / open-source computational
tooling.

\begin{enumerate}

\item \textbf{FLASH-regime extrapolation.}
Does the UNIVERSE repair-kinetics sub-model extrapolate correctly into the
FLASH dose-rate regime ($>40$~Gy/s), where transient DSB density and
inter-track saturation should differ qualitatively from the 6--53~Gy/min
HIT SOBP rates the paper actually tests?

\emph{Basis.} The paper's dose-rate coverage tops out at $\sim$1~Gy/s
(53~Gy/min); FLASH-relevant rates are 3--4 orders of magnitude higher. The
Poisson-per-slab domain-collision model may saturate at very high rates in a
way not tested against data. Fig.~4 left panel plateaus above 100~Gy/min in
the model but no measurement anchors that plateau.

\emph{Next steps.} (i) Sweep the $R_{TD50}$ curve from 3.75~Gy/min to
$10^{5}$~Gy/min at fixed $N_t$ and check the plateau value.
(ii) Repeat with $N_t$ rescaled so the per-slab dose stays below the
intra-domain collision threshold.
(iii) Compare against FLASH radiobiology datasets (Favaudon 2014;
Montay-Gruel 2019) after mapping their surrogate endpoints onto the
UNIVERSE survival framework.

\item \textbf{Pathway-explicit repair kinetics.}
How does the two-class isolated/complex DSB model compare against a modern
pathway-explicit repair kinetics model (NHEJ fast + HR slow + alt-EJ)
parameterized from ionizing-radiation-response gene-expression data?

\emph{Basis.} UNIVERSE lumps all repair into two effective rate constants
($T_{iDSB}^{1/2}$, $T_{cDSB}^{1/2}$). Modern DSB repair biology tracks
distinct pathways with cell-cycle-dependent choice. The paper's
cell-line-specific $T_{iDSB}^{1/2}$ range of 4--11~min plausibly aggregates
20-min fast NHEJ + hours-scale HR into one exponential.

\emph{Next steps.} (i) Replace the two-rate exponential lifetime with a
mixture of NHEJ+HR exponentials with published rate constants
(Mladenov 2016).
(ii) Fit $K_{iDSB}$, $K_{cDSB}$, and pathway mixing fractions jointly
against the same RSC TD50 tables.
(iii) Test whether an explicit-pathway model narrows or widens the
residuals against the with-repair fit in the paper's Table~1.

\item \textbf{Transfer to clinical tumor cell lines.}
Do the paper's DU145 and RSC $K_{iDSB}$ and $K_{cDSB}$ parameters transfer
to specific human tumor cell lines used in modern ion-beam clinical
planning (H460, A549, U87, chordoma primaries), or must they be refit per
line?

\emph{Basis.} The paper reports only DU145 (in vitro, prostate) and RSC
(in vivo, rat). Clinical carbon and proton planning at HIT/MedAustron uses
cell-line-agnostic $\alpha/\beta$ ratios. Whether the $K$ and $T$
parameters need per-line refitting or transfer via $\alpha/\beta$ scaling is
not addressed by the paper.

\emph{Next steps.} (i) Extract published DSB repair kinetics (via
$\gamma$-H2AX foci decay) for H460, A549, U87, chordoma lines from the DDR
literature. (ii) Fit the four kinetic parameters to those foci curves.
(iii) Predict RBE for each line at HIT SOBP LET and compare against
published in-vitro clonogenic RBE data.

\item \textbf{Joint parameter uncertainty and Bayesian fit.}
How sensitive is the reproduced $R_{TD50}$ to joint uncertainty in
$\{K_{iDSB},\ K_{cDSB},\ T_{iDSB}^{1/2},\ T_{cDSB}^{1/2}\}$, and does the
paper's point estimate lie within a 68\% credible ellipse from a proper
Bayesian fit to Karger 2003 + Karger 2006 + Saager 2018 + Hintz 2022?

\emph{Basis.} The paper reports parameter point estimates only; no
posterior, no covariance, no goodness-of-fit $\chi^{2}$. Our 0.83\% MAD
agreement is at the MC noise floor, so we cannot distinguish the paper's
fit from nearby fits without joint parameter uncertainty quantification.

\emph{Next steps.} (i) Wrap the MC survival model in a likelihood function
against the four TD50 datasets. (ii) Run emcee or dynesty over the 4D
parameter space with weakly-informative priors. (iii) Report 68\% and 95\%
marginal intervals and check whether the paper's point estimate is within
the 1-$\sigma$ ellipse or is a MAP under a specific implicit prior.

\item \textbf{ML-driven parameter identification.}
Can neural or symbolic-regression-based ML identify a lower-dimensional
latent repair-kinetics parameterization that matches the two-rate UNIVERSE
fit quality across all tested cell lines with fewer free constants?

\emph{Basis.} The UNIVERSE two-rate model has 4 free parameters per cell
line. It is plausible that a 2- or 3-parameter surrogate learned across
cell lines (perhaps parameterized by cell-cycle fraction and BRCA1
competence) reproduces the same $R_{TD50}$ predictions with tighter
cross-line consistency, revealing biology hidden by per-line refitting.

\emph{Next steps.} (i) Assemble a corpus of (cell-line, $\alpha$, $\beta$,
$K_{iDSB}$, $K_{cDSB}$, $T_{iDSB}^{1/2}$, $T_{cDSB}^{1/2}$) tuples across
published UNIVERSE / LEM-IV fits. (ii) Fit a low-rank factorization or
symbolic regressor mapping bio-features (cell-cycle distribution,
DDR-gene expression) to the four repair constants. (iii) Compare held-out
cell-line $R_{TD50}$ predictions between the ML surrogate and per-line
refitting.

\end{enumerate}
